\chapter{Ultrapowers with $\I$: Construction and Properties}
\label{ch:ultrapowers}

\section{Introduction}

Ultrapower constructions provide a powerful tool for building models of extended theories. In this chapter, we develop the theory of ultrapowers that accommodate the indeterminate element $\I$, establishing both the construction methods and the preservation properties that make ultrapowers fundamental to our framework.

\section{Ultrafilters: Review and Extension}

\subsection{Classical Ultrafilters}

\begin{definition}[Ultrafilter]
An \textbf{ultrafilter} $\mathcal{U}$ on a set $I$ is a collection of subsets of $I$ satisfying:
\begin{enumerate}
    \item $I \in \mathcal{U}$ and $\emptyset \notin \mathcal{U}$ (non-triviality)
    \item If $A, B \in \mathcal{U}$, then $A \cap B \in \mathcal{U}$ (closure under intersection)
    \item If $A \in \mathcal{U}$ and $A \subseteq B \subseteq I$, then $B \in \mathcal{U}$ (upward closure)
    \item For all $A \subseteq I$, either $A \in \mathcal{U}$ or $I \setminus A \in \mathcal{U}$ (ultraness)
\end{enumerate}
\end{definition}

\begin{definition}[Principal vs. Non-Principal]
An ultrafilter $\mathcal{U}$ on $I$ is:
\begin{itemize}
    \item \textbf{Principal} if $\mathcal{U} = \{A \subseteq I : i_0 \in A\}$ for some $i_0 \in I$
    \item \textbf{Non-principal} if it contains no finite sets
\end{itemize}
\end{definition}

\subsection{Ultrafilters on $\N$ and Indeterminacy}

For our purposes, we focus on non-principal ultrafilters on $\N$:

\begin{proposition}[Existence of Non-Principal Ultrafilters]
\label{prop:nonprincipal-exist}
Assuming the Axiom of Choice, there exist non-principal ultrafilters on $\N$.
\end{proposition}

\begin{proof}[Proof Sketch]
Use Zorn's lemma on the poset of filters containing all cofinite sets, ordered by inclusion. Any maximal element is a non-principal ultrafilter.
\end{proof}

\begin{remark}
The existence of non-principal ultrafilters is equivalent to a weak form of AC and cannot be proven in ZF alone.
\end{remark}

\section{Ultrapower Construction}

\subsection{Classical Ultrapower}

\begin{definition}[Ultrapower]
Let $\mathcal{M}$ be a structure with domain $M$ and let $\mathcal{U}$ be an ultrafilter on $I$. The \textbf{ultrapower} $\mathcal{M}^I/\mathcal{U}$ is defined as follows:
\begin{itemize}
    \item Domain: Equivalence classes of functions $f : I \to M$ under the relation:
    \[
    f \sim_{\mathcal{U}} g \iff \{i \in I : f(i) = g(i)\} \in \mathcal{U}
    \]
    \item Operations and relations: Applied pointwise and then modulo $\mathcal{U}$
\end{itemize}
\end{definition}

\subsection{Extended Ultrapower with $\I$}

We now extend the ultrapower construction to accommodate $\I$:

\begin{definition}[Extended Ultrapower $M^*$]
\label{def:extended-ultrapower}
Let $\mathcal{U}$ be a non-principal ultrafilter on $\N$. Define:
\[
M^* = \prod_{n \in \N} M / \mathcal{U} \cup \{\I\}
\]
where sequences $(a_n)$ satisfying certain indeterminacy conditions are identified with $\I$.
\end{definition}

Specifically, we identify a sequence $(a_n)$ with $\I$ if:
\[
\lim_{n \to \mathcal{U}} (a_n - a_n) \neq 0
\]
in the sense that the set $\{n : a_n - a_n \text{ is undefined or divergent}\} \in \mathcal{U}$.

\subsection{Properties of the Extended Ultrapower}

\begin{theorem}[Well-Definedness of $M^*$]
The extended ultrapower $M^*$ is well-defined, i.e., the equivalence relation is compatible with the operations.
\end{theorem}

\begin{proof}
We verify that if $[f]_{\mathcal{U}} = [f']_{\mathcal{U}}$ and $[g]_{\mathcal{U}} = [g']_{\mathcal{U}}$, then:
\[
[f \oplus g]_{\mathcal{U}} = [f' \oplus g']_{\mathcal{U}}
\]

Let $A = \{n : f(n) = f'(n)\} \in \mathcal{U}$ and $B = \{n : g(n) = g'(n)\} \in \mathcal{U}$.

Then $A \cap B \in \mathcal{U}$ by closure, and for all $n \in A \cap B$:
\[
(f \oplus g)(n) = f(n) \oplus g(n) = f'(n) \oplus g'(n) = (f' \oplus g')(n)
\]

Thus $\{n : (f \oplus g)(n) = (f' \oplus g')(n)\} \supseteq A \cap B \in \mathcal{U}$, so $[f \oplus g]_{\mathcal{U}} = [f' \oplus g']_{\mathcal{U}}$.

The case involving $\I$ requires checking that indeterminate sequences are mapped to $\I$ consistently under operations, which follows from Axiom $\I_4$ (absorption).
\end{proof}

\section{Łoś's Theorem for Extended Structures}

\subsection{Classical Łoś's Theorem}

\begin{theorem}[Łoś's Theorem -- Classical]
\label{thm:los-classical}
For any first-order formula $\phi(x_1, \ldots, x_n)$ and functions $f_1, \ldots, f_n : I \to M$:
\[
\mathcal{M}^I/\mathcal{U} \models \phi([f_1]_{\mathcal{U}}, \ldots, [f_n]_{\mathcal{U}}) \iff \{i \in I : \mathcal{M} \models \phi(f_1(i), \ldots, f_n(i))\} \in \mathcal{U}
\]
\end{theorem}

\subsection{Modified Łoś's Theorem for $\I$}

When $\I$ is present, Łoś's theorem requires modification:

\begin{theorem}[Łoś's Theorem with $\I$]
\label{thm:los-extended}
For first-order formulas $\phi$ and functions $f_1, \ldots, f_n : I \to M \cup \{\I\}$:
\[
\mathcal{M}^*/\mathcal{U} \models \phi([f_1], \ldots, [f_n]) \iff \{i : \mathcal{M} \models \phi(f_1(i), \ldots, f_n(i))\} \in \mathcal{U}
\]
with the proviso that formulas involving $\I$ are evaluated using $\I$-logic (three-valued) rather than classical logic.
\end{theorem}

\begin{proof}[Proof Sketch]
The proof proceeds by induction on formula complexity:

\textbf{Base case}: For atomic formulas $P(x)$:
\begin{itemize}
    \item If $[f]_{\mathcal{U}} = c \in M$, standard Łoś applies
    \item If $[f]_{\mathcal{U}} = \I$, then $P(\I)$ has indeterminate truth value by definition
\end{itemize}

\textbf{Inductive cases}:
\begin{itemize}
    \item \textbf{Conjunction}: $\phi \land \psi$ uses the $\I$-logic truth table for $\land$
    \item \textbf{Negation}: $\neg \phi$ maps $\I$-values to $\I$-values
    \item \textbf{Quantifiers}: $\exists x \, \phi(x)$ is true if witnessed by either a determinate value or $\I$
\end{itemize}

The key modification is that truth values can be indeterminate, and the ultrafilter determines which indeterminacies are "large" (in $\mathcal{U}$).
\end{proof}

\section{Preservation of Axioms}

\subsection{Axiom Preservation in Ultrapowers}

\begin{theorem}[Preservation of $\I$-Axioms]
\label{thm:axiom-preservation}
Let $\mathcal{M}$ be a model of $\text{ZF} + \text{AC} + \I_1 + \cdots + \I_{15}$. Then the ultrapower $\mathcal{M}^*/\mathcal{U}$ also satisfies all these axioms.
\end{theorem}

\begin{proof}
We verify each axiom class:

\textbf{ZF + AC}: These are first-order and preserved by classical Łoś.

\textbf{Axiom $\I_1$}: For any cardinal $[\kappa]_{\mathcal{U}}$ in the ultrapower:
\begin{align*}
[\kappa]_{\mathcal{U}} \ominus [\kappa]_{\mathcal{U}} &= [\kappa \ominus \kappa]_{\mathcal{U}} \quad \text{(operations defined pointwise)} \\
&= [\I]_{\mathcal{U}} \quad \text{(by $\I_1$ in $\mathcal{M}$)} \\
&= \I \quad \text{(constant functions represent themselves)}
\end{align*}

\textbf{Axiom $\I_4$ (Absorption)}: For any $[\kappa]_{\mathcal{U}}$:
\begin{align*}
\I \oplus [\kappa]_{\mathcal{U}} &= [\I]_{\mathcal{U}} \oplus [\kappa]_{\mathcal{U}} \\
&= [\I \oplus \kappa]_{\mathcal{U}} \\
&= [\I]_{\mathcal{U}} \quad \text{(by $\I_4$ in $\mathcal{M}$)} \\
&= \I
\end{align*}

\textbf{Axiom $\I_7$ (Non-Triviality)}: Since $\mathcal{M} \models \I \neq 0$, by Łoś:
\[
\{i : \I \neq 0\} = I \in \mathcal{U}
\]
Thus $\mathcal{M}^*/\mathcal{U} \models \I \neq 0$.

Similar arguments apply to the remaining axioms.
\end{proof}

\section{Saturation and Transfer}

\subsection{Saturation Properties}

\begin{definition}[$\kappa$-Saturation]
A structure $\mathcal{M}$ is \textbf{$\kappa$-saturated} if for every set $\Phi(x)$ of fewer than $\kappa$ formulas with parameters from $M$, if every finite subset of $\Phi$ is satisfiable in $\mathcal{M}$, then $\Phi$ is satisfiable in $\mathcal{M}$.
\end{definition}

\begin{theorem}[Saturation of Ultrapowers]
If $|M| < |\mathcal{U}|$, then $\mathcal{M}^I/\mathcal{U}$ is $|M|^+$-saturated.
\end{theorem}

\subsection{Transfer Principles}

\begin{theorem}[Partial Transfer for $\I$]
First-order properties involving only determinate elements transfer between $\mathcal{M}$ and $\mathcal{M}^*/\mathcal{U}$. Properties essentially involving $\I$ may not transfer without additional conditions.
\end{theorem}

\section{Applications to Consistency Proofs}

\subsection{Relative Consistency via Ultrapowers}

\begin{theorem}[Consistency of Extended Theory]
\label{thm:ultrapower-consistency}
If there exists a model $\mathcal{M}$ of ZF + AC, then there exists a model $\mathcal{M}^*$ of ZF + AC + $\I_1$--$\I_{15}$.
\end{theorem}

\begin{proof}
Start with $\mathcal{M} \models \text{ZF} + \text{AC}$. Construct:
\begin{enumerate}
    \item Choose a non-principal ultrafilter $\mathcal{U}$ on $\N$
    \item Form the extended ultrapower $\mathcal{M}^* = \mathcal{M}^\N/\mathcal{U} \cup \{\I\}$
    \item Define operations on $\mathcal{M}^*$ extending those of $\mathcal{M}$
    \item Verify Axioms $\I_1$--$\I_{15}$ using Theorem~\ref{thm:axiom-preservation}
\end{enumerate}

By the construction, $\mathcal{M}^* \models \text{ZF} + \text{AC} + \I_1 + \cdots + \I_{15}$.
\end{proof}

\section{Examples and Computations}

\subsection{Ultrapower of $\R$}

\begin{example}[Extended Hyperreals]
Consider $\mathcal{R} = (\R, +, \cdot, <)$. The ultrapower $\mathcal{R}^* = \R^\N/\mathcal{U}$ contains:
\begin{itemize}
    \item Standard reals: Constant sequences $[c, c, c, \ldots]$
    \item Infinitesimals: $[\frac{1}{1}, \frac{1}{2}, \frac{1}{3}, \ldots]$
    \item Infinities: $[1, 2, 3, \ldots]$
    \item Indeterminate: $\I$ from $[0, 1, 0, 1, \ldots]$ (oscillating)
\end{itemize}
\end{example}

\subsection{Cardinal Ultrapowers}

\begin{example}[Cardinal Ultrapower]
Let $\mathcal{C} = (\text{Card}, +, \cdot, \aleph_0, \aleph_1, \ldots)$. In $\mathcal{C}^*/\mathcal{U}$:
\begin{itemize}
    \item $[\aleph_0] \ominus [\aleph_0] = \I$
    \item $[\aleph_1] \ominus [\aleph_1] = \I$
    \item All self-subtractions yield the same $\I$, confirming uniqueness
\end{itemize}
\end{example}

\section{Connection to Nonstandard Analysis}

The extended ultrapower construction generalizes Robinson's nonstandard analysis:

\begin{center}
\begin{tabular}{|l|l|}
\hline
\textbf{Robinson's NSA} & \textbf{Our Framework} \\
\hline
Infinitesimals & Indeterminate elements \\
Transfer principle & Modified Łoś theorem \\
${}^*\R \supseteq \R$ & $\KI \supseteq \R$ \\
$\epsilon \approx 0$ & $\I$ (fundamentally distinct from 0) \\
\hline
\end{tabular}
\end{center}

\section{Computational Implementation}

\subsection{Python Simulation}

\begin{verbatim}
class Ultrafilter:
    """Non-principal ultrafilter on N (approximation)."""
    
    def contains(self, subset):
        """Check if subset is in ultrafilter (finite approximation)."""
        # Use cofinite sets as approximation
        return len([n for n in range(1000) if n not in subset]) < 100

class Ultrapower:
    """Ultrapower construction M^N/U."""
    
    def __init__(self, base_structure, ultrafilter):
        self.base = base_structure
        self.U = ultrafilter
    
    def equivalence_class(self, sequence):
        """Form equivalence class [f]_U."""
        # Check for indeterminacy
        if self.is_indeterminate(sequence):
            return IndeterminateElement()
        return EquivalenceClass(sequence, self.U)
    
    def is_indeterminate(self, seq):
        """Check if sequence represents I."""
        # Simplified: check for oscillation or self-subtraction
        diffs = [seq[i] - seq[i] for i in range(len(seq))]
        return not all(d == 0 for d in diffs)
\end{verbatim}

\section{Summary}

This chapter has established:

\begin{enumerate}
    \item Extended ultrapower constructions accommodating $\I$
    \item Modified Łoś's theorem for three-valued logic
    \item Preservation of $\I$-axioms under ultrapower
    \item Applications to consistency proofs
    \item Connections to nonstandard analysis
\end{enumerate}

Ultrapowers provide the model-theoretic foundation for studying $\I$ in extended structures, enabling rigorous analysis of indeterminacy in mathematical systems.

\section{Exercises}

\begin{enumerate}
    \item Prove that non-principal ultrafilters on $\N$ contain all cofinite sets.
    \item Construct an explicit non-principal ultrafilter using Zorn's lemma.
    \item Verify the modified Łoś theorem for quantified formulas.
    \item Show that the ultrapower $\R^\N/\mathcal{U}$ is $\aleph_1$-saturated.
    \item Implement the ultrapower construction in Lean 4 and prove Theorem~\ref{thm:axiom-preservation}.
\end{enumerate}
